\documentclass[10pt,a4paper]{amsart}
\usepackage[margin=19mm]{geometry}
\usepackage{amsmath,amssymb,mathtools}
\usepackage[scr=rsfs,cal=euler]{mathalfa}
\usepackage{xfrac}
\usepackage[unicode,hidelinks,bookmarks=false]{hyperref}
\newcommand{\ie}{i.e.\ }
\newcommand{\cf}{cf.\ }
\newcommand{\resp}{respectively\ }
\newcommand{\Subsection}{\subsection}
\providecommand{\nobreakdash}{\nobreak\hbox{-}}
\setlength{\emergencystretch}{2em}
\multlinegap0pt
\usepackage{mathtools}
\newtheorem{theorem}{Theorem}[section]
\newtheorem{proposition}[theorem]{Proposition}
\newtheorem{lemma}[theorem]{Lemma}
\newtheorem{corollary}[theorem]{Corollary}
\newtheorem{remark}[theorem]{Remark}
\newcommand{\E}{\mathbb E}
\newcommand{\R}{\mathbb R}
\newcommand{\Prob}{\mathbb P}
\newcommand{\Cov}{\operatorname{Cov}}
\newcommand{\Var}{\operatorname{Var}}
\newcommand{\arctanh}{\operatorname{arctanh}}
\hypersetup{pdftitle={Critical and near-critical influence bounds for ferromagnetic Ising models},pdfauthor={Yan Ru Pei}}
\begin{document}
\title[Critical and near-critical influence bounds for ferromagneti]{Critical and near-critical influence bounds for ferromagnetic Ising models}
\author{Yan Ru Pei}
\date{}
\maketitle
\noindent\emph{Source and licence.} Critical and near-critical influence bounds for ferromagnetic Ising models. Version arXiv:2609.20866v1 (CC BY 4.0). This material is distributed under the Creative Commons Attribution 4.0 International licence, http://creativecommons.org/licenses/by/4.0/, as recorded in the arXiv OAI licence metadata for this exact version and shown on the abstract page https://arxiv.org/abs/2609.20866v1. Full rights notice: licenses/arxiv-2609.20866-CC-BY-4.0.txt.\par\smallskip
\noindent\textit{Editorial scope.} Counted unit: the original \texttt{lemma} environment \texttt{lem:cavity} at source lines 155--168 together with its complete proof at lines 169--184; the delivered document keeps source lines 147--185 so that every referenced label and every citation resolves inside it. Native editable source is the paper's own concatenated TeX; the AMS wrapper is editorial formatting only. No publisher class, upstream script or unrelated artwork is redistributed.
\section{A pointwise cavity estimate}

We first work in a uniform field $h>0$. Denote the magnetization at $u$
in a graph $H$ by $m_u^H$, and suppress $h$ in the notation.
The Griffiths correlation inequalities imply that spin-product
expectations increase when nonnegative edges are added.
These are the classical inequalities used in~\cite[Section~4]{CDKP}.


\begin{lemma}\label{lem:cavity}
Suppose every edge has coupling $\beta$, and let
$t=\tanh\beta$ and $f_\beta(x)=\arctanh(t\tanh x)$.
For each oriented edge put
$L_{u\to v}=\arctanh(m_u^{G-uv})$. Then
\begin{align}
 L_{u\to v}
    &\le h+\sum_{w\in N(u)\setminus\{v\}}f_\beta(L_{w\to u}),
       \label{eq:cavity}\\
 \arctanh(m_u^G)
    &\le h+\sum_{w\in N(u)}f_\beta(L_{w\to u}).
       \label{eq:root}
\end{align}
\end{lemma}

\begin{proof}
For an edge $uv$ of $H$, let
$A=m_u^{H-uv}$, $B=m_v^{H-uv}$, and
$C=\E_{H-uv,h}[\sigma_u\sigma_v]$.
Adding the edge gives the exact identity
$m_u^H=(A+tB)/(1+tC)$.
Since $C\ge AB$ and $A+tB\ge0$,
\[
 \arctanh(m_u^H)\le\arctanh A+\arctanh(tB).
\]
Remove the edges incident to $u$ one at a time.
At the removal of $uw$, monotonicity bounds the neighbor
magnetization by $m_w^{G-uw}$.
The final isolated vertex has magnetization $\tanh h$.
Starting with $G-uv$ or $G$ proves the two inequalities.
\end{proof}

\begin{thebibliography}{99}
\bibitem[CDKP]{CDKP} C. Carlson, E. Davies, A. Kolla and W. Perkins, \emph{Computational thresholds for the fixed-magnetization Ising model}, \href{https://arxiv.org/abs/2111.03033}{arXiv:2111.03033}, 2021.
\end{thebibliography}
\end{document}
